\subsection{The ordinary-type topology on a curve}
\label{subsec:curve-ordinary-topology}

The topology of \cref{sec:ordinary-type-topology} also uses the ordinary envelope, but it asks for total variation only of the atomic curvature. On ideal models that is the whole polar measure, so it recovers the complete metric topology above. On arbitrary ordinary models, diffuse polar mass can concentrate without being seen by this variation coordinate.

Retain $Y,\theta,A$ and the decomposition $\mu_u=D_u+\sigma_u$ above. Use $\eta=\theta/A$ for $L^1$ norms. For ordinary classes put
\begin{equation}\label{eq:curve-ordinary-type-distance}
 d^{\mathrm{ot}}_\theta([u],[v])
 =\|P_\theta[u]-P_\theta[v]\|_{L^1(\eta)}
     +\|D_u-D_v\|_{\mathrm{TV}}.
\end{equation}

\begin{theorem}\label{thm:curve-ordinary-type-topology}\index{Singularity type!curve pseudometric}
The finite pseudometric $d^{\mathrm{ot}}_\theta$ induces $\tau_{\theta,0}$ on the whole ordinary-type space. For $0\leq s\leq A$, the canonical ideal classes
\[
 \mathcal C_s=\{[q_D]:D\text{ positive atomic},\ D(Y)\leq s\}
\]
form a complete metric space both for $d^{\mathrm{ot}}_\theta$ and for $\|D-E\|_{\mathrm{TV}}$. These metrics induce the same topology. The space $\mathcal C_s$ is not totally bounded if $s>0$, and $\mathcal C_0=\{[0]\}$.
\end{theorem}
\begin{proof}
A smooth compact modification of a smooth curve is an isomorphism. Its divisorial distance is exactly atomic total variation, since the divisor weight in dimension one is one. Thus the two summands of \eqref{eq:curve-ordinary-type-distance} are defining coordinates of $\tau_{\theta,0}$.

Conversely, suppose a net converges for $d^{\mathrm{ot}}_\theta$. The ideal envelope of $u$ is $q_{D_u}$ by \cref{lem:curve-atomic-ideal-model}. Atomic total variation implies $L^1$ convergence of these envelopes by \cref{thm:curve-model-capacity-convergence}. This theorem also gives continuity for nets: its domain and target are metric spaces, so its sequential assertion is equivalent to continuity. It remains to check the ideal relations.

At a point $p$ with local coordinate $z$, the curve ideal formula is
\[
 \mathcal I(tu)_p=(z^{\lfloor tD_u(\{p\})\rfloor}).
\]
This follows from \cref{thm:local-i-inputs,lem:smooth-center-multiplier-ideal}: on a smooth curve the point Lelong number determines the ideal, and the logarithmic germ with that number has the displayed ideal. Write $D_j=D_{u_j}$, $D=D_u$ and fix $t>0$, $a>1$. The set $F=\{p:D(\{p\})\geq1/(2t)\}$ is finite. At its points, $aD_j(\{p\})\geq D(\{p\})$ and $aD(\{p\})\geq D_j(\{p\})$ eventually. Outside $F$, the bound
\[
 \sup_p|D_j(\{p\})-D(\{p\})|\leq\|D_j-D\|_{\mathrm{TV}}\longrightarrow0
\]
makes both $\mathcal I(tu)_p$ and $\mathcal I(tu_j)_p$ unit ideals for large $j$. The two desired inclusions therefore hold everywhere. This proves the asserted topology.

The map $D\mapsto q_D$ is injective. Positive atomic measures of mass at most $s$ are closed in total variation: a limit is positive, has this mass bound, and is supported on the countable union of the supports of a converging sequence. Completeness of finite measures gives completeness for total variation. A $d^{\mathrm{ot}}_\theta$-Cauchy sequence is total-variation Cauchy; the resulting limit remains atomic, and \cref{thm:curve-model-capacity-convergence} supplies convergence of the envelope coordinate. This proves its completeness as well as equality of the two topologies.

If $s>0$, fix $0<b\leq s$ and infinitely many distinct points $p_j$. The measures $b\delta_{p_j}$ have pairwise total-variation distance $2b$, proving failure of total boundedness for both metrics. At $s=0$ only the zero measure remains.
\end{proof}

By \cref{lem:curve-atomic-ideal-model}, $\mathcal C_{A-\delta}$ is exactly the canonical ideal locus of mass at least $\delta$, for $0\leq\delta\leq A$. The entire canonical locus $\mathcal C_A$ is complete, including zero mass. Its positive-mass part is incomplete: $D_j=A(1-1/j)\delta_p$ is Cauchy and converges only to the zero-mass boundary $A\delta_p$.

The same statement on actual ordinary classes is a pseudometric assertion. More precisely, put
\[
 \mathcal G_{\theta,\delta}
 =\left\{[u]:P_\theta[u]=P_\theta[u]_{\mathcal I},\quad
       \int_Y\mathrm{MA}_\theta(P_\theta[u]_{\mathcal I})\geq\delta\right\}.
\]
For $0\leq\delta\leq A$ this is complete for $d^{\mathrm{ot}}_\theta$, since every point is at distance zero from $[q_{D_u}]\in\mathcal C_{A-\delta}$. Limits need not be unique as ordinary classes. It is not totally bounded when $\delta<A$; at $\delta=A$ the pseudometric vanishes identically, even though \cref{prop:ordinary-topology-nonseparation} supplies distinct classes. Its strictly positive-mass part without a uniform lower bound is incomplete by the same one-atom sequence.

\subsection{Concentration of diffuse polar curvature}

Atomic variation is indispensable in \cref{thm:curve-ordinary-type-topology}. In the example of \cref{prop:positive-mass-ideal-discontinuity}, choose the rational approximations so that the largest atom of $\nu_j$ tends to zero. The ordinary and ideal envelopes both tend to zero in $L^1$, and for each $t,a$ all the controlled ideal relations with zero eventually hold. Their currents converge weakly to $\theta$. Nevertheless,
\[
 d^{\mathrm{ot}}_\theta([r_j],[0])
 \geq\|D_{r_j}\|_{\mathrm{TV}}=1-\tau>0.
\]
Thus these weaker conditions do not imply the ordinary-type topology even at a fixed positive canonical mass.

The converse difficulty concerns diffuse polar mass: the pseudometric controls its envelope in $L^1$, but not the total variation of the diffuse measure itself.

\begin{proposition}\label{prop:curve-ordinary-topology-incomplete}\index{Singularity type!incomplete ordinary space}
For the probability Fubini--Study form $\theta$ on $\mathbb P^1$, the whole space $(\mathfrak S(\mathbb P^1,\theta),d^{\mathrm{ot}}_\theta)$ is incomplete. The failure occurs for ordinary models of one fixed positive non-pluripolar mass, all with trivial multiplier ideals.
\end{proposition}
\begin{proof}
Choose a nonatomic polar probability measure $\sigma$ with compact support in an affine chart, using the Cantor construction in \cref{ex:diffuse-polar-zero-model}. Fix $0<c<1$ and a point $p$ in that chart. Translate and shrink the support towards $p$, obtaining nonatomic polar probabilities $\sigma_j$ supported in discs of radii tending to zero about $p$. Let $p_j=p_{c\sigma_j}$ be the ordinary models of \cref{lem:curve-polar-model}. Their polar measures are $c\sigma_j$, their masses are $1-c$, and $D_{p_j}=0$. Their ideal envelopes are zero.

We claim $p_j\to q_{c\delta_p}$ in $L^1$. Write $q_y=q_{c\delta_y}$. Rotations preserving $\theta$ can be chosen continuously near $p$ to send $p$ to $y$. They transport $q_p$ to $q_y$, and continuity of translations in $L^1$ gives $q_y\to q_p$ there as $y\to p$. Set
\[
 w_j=\int_{\mathbb P^1}q_y\,\mathrm d\sigma_j(y).
\]
The common normalized $L^1$ bound makes this integral well-defined. It is nonpositive, its curvature dominates $c\sigma_j$, and $w_j\to q_p$ in $L^1$. The envelope characterization \eqref{eq:curve-polar-domination} therefore gives $w_j\leq p_j$.

Normalized compactness supplies an $L^1$ cluster limit $v$ along every subsequence. The inequalities $\theta+\ddc p_j\geq c\sigma_j$ pass to $\theta+\ddc v\geq c\delta_p$. Thus $v\leq q_p$ by the ideal-envelope characterization. The lower bound from $w_j$ gives $v\geq q_p$ almost everywhere, hence everywhere. This identifies every cluster limit and proves the claim.

Since $P_\theta[p_j]=p_j$ and $D_{p_j}=0$, the classes $[p_j]$ are $d^{\mathrm{ot}}_\theta$-Cauchy. A limit $[v]$ would have $D_v=0$ and $P_\theta[v]=q_p$. This is impossible: ordinary envelopes preserve all multiplier ideals and hence all atoms, whereas $q_p$ has atom $c$ at $p$. The sequence therefore has no limit in the ordinary-type space.
\end{proof}

The two topologies on ordinary models are also different. Choose the Cantor support of $\sigma$ inside a positive real interval in the affine chart, away from zero. The rotations $R_j(z)=\mathrm e^{\mathrm i\varepsilon_j}z$, with $0<\varepsilon_j\to0$, move this support off itself. Put $p=p_{c\sigma}$ and $\widetilde p_j=p\circ R_j^{-1}$. Rotation invariance of the envelope gives ordinary models with polar measures $c(R_j)_*\sigma$ and no atoms. Continuity of rotations in $L^1$ yields $d^{\mathrm{ot}}_\theta([\widetilde p_j],[p])\to0$, whereas disjoint polar supports give $d_\theta(\widetilde p_j,p)=2c$ by \cref{thm:curve-model-volume-distance}. Conversely, polar total variation controls both terms of \eqref{eq:curve-ordinary-type-distance}, by \cref{thm:curve-model-capacity-convergence}. Thus $\tau_{\theta,0}$ is strictly weaker on ordinary models, although its restriction to ideal models agrees with their complete metric topology. The analytic-closure theorem alone does not extend the positive-volume model topology.
